\documentclass[12pt,a4paper]{article}

% PACOTES
\usepackage[portuguese]{babel}
\usepackage[utf8]{inputenc}
\usepackage[T1]{fontenc}
\usepackage{makeidx}
\usepackage{xspace}
\usepackage{graphicx,color,times}
\usepackage{fancyhdr}
\usepackage{latexsym}
\usepackage[printonlyused]{acronym} % Pacote acronym aqui
\usepackage{float}
\usepackage{listings}
% remove duplicate "index"
\usepackage[nottoc]{tocbibind}
\usepackage{natbib}
\usepackage{hyperref}
\usepackage{xcolor}
\usepackage{fix-cm}
\usepackage{fourier}
\usepackage[scaled=.92]{helvet}
% show references names
\usepackage{nameref}

% CONFIGURAÇÕES HYPERREF
\hypersetup{
    colorlinks=true,            
    linkcolor=black,            
    citecolor=blue,             
    urlcolor=blue,              
    filecolor=blue,             
    breaklinks=true,            
    pdftitle={A Governança da Internet},   
    pdfauthor={Tomás Simões Duarte},           
    pdfsubject=
    {
        As decisões no design e desenvolvimento 
        de tecnologias que são necessárias para manter a Internet operacional 
        com a implementação
        de políticas novas em volta das mesmas
    }, 
    pdfkeywords={Governança, Internet},   
}

% DEFINIÇÃO DE ACRÔNIMOS
\begin{document}

% Frontmatter
\thispagestyle{empty}
\setcounter{page}{-1}

\begin{center}
  \begin{Huge}
    \textbf{Universidade da Beira Interior}
  \end{Huge}
\end{center}

\begin{center}
  \begin{Huge}
    Departamento de Informática
  \end{Huge}
\end{center}

\vspace{0,07cm}
\begin{figure}[!htb]
  \centering
  \includegraphics[width=191pt]{images/ubi-fe-di.png}
\end{figure}

\vspace{0.5cm}
\begin{center}
  \begin{Large}
    \textbf{A Governancia da Internet}
  \end{Large}
\end{center}


\vspace{0.5cm}
\begin{center}
  \begin{normalsize}
    \begin{large}
      Elaborado por:
    \end{large}
  \end{normalsize}
\end{center}

\vspace{0.2cm}
\begin{center}
  \begin{large}
    \textbf{Tomás Simões [52585]}
  \end{large}
\end{center}

\vspace{0,5cm}
\begin{center}
  \begin{normalsize}
    \begin{large}
      Orientador:
    \end{large}
  \end{normalsize}
\end{center}

\vspace{0.2cm}
\begin{center}
  \begin{large}
    \textbf{Professor Doutor Bruno Miguel Silva}
  \end{large}
\end{center}



\vspace{0.5cm}
\begin{center}
  \begin{normalsize}
    \today
  \end{normalsize}
\end{center}

\let\cleardoublepage\clearpage
\tableofcontents

% Mainmatter
\acresetall
\section{Introdução}
\label{chap:intro}
Este relatório, com o tema de "Governance da Internet" foi feito no contexto da unidade curricular de Redes de Computador, no projeto
enquanto frequentava a \ac{UBI}.

Atualmente vivemos em uma sociedade que é governada pela tecnologica, apesar de a maioria dos individuos
não saberem quem a governa.

Ao longo dos anos tivemos diferentes perpetivas sobre o uso desta,
desde o uso militar até simplemente podermos usar o Youtube num frigorifico. Por essa razão
é necessário atualizar e dialogar sobre como deviamos nos comportar face a questões políticas, sociais e filosoficas da rede de redes
de modo a termos um espaço seguro para estar.

Este documento fala sobre o aparecimento do termo "Governancia da Internet", um campo de estudo e pesquisa,
complexo e com diversas intrepetações sobre quem  dentro da Internet.

\subsection{Organização do Documento}
\label{sec:organ}

De modo a refletir o trabalho que foi feito, este documento encontra-se estruturado da seguinte forma:
\begin{enumerate}
  \item O 1º capítulo -- \textbf{Introdução} -- apresenta o projeto, a motivação para a sua escolha, o enquadramento para o mesmo, os seus objetivos e a respetiva organização do documento.
  \item O 2º capítulo -- \textbf{Tecnologias Utilizadas} -- descreve os conceitos mais importantes no âmbito deste projeto, bem como as tecnologias utilizadas durante do desenvolvimento da aplicação.
  \item ...
\end{enumerate}

\section{O que é a Governancia da Internet?}
\label{sec:internetgovernance}
A Governancia da Internet é definida como "decisões no \textit{design} e desenvolvimento 
de tecnologias que são necessárias para manter a Internet operacional com a implementação
de políticas novas em volta das mesmas".

Ao longo do tempo, pesquisadores assim como profissionais na área têm tentado expandir o 
conhecimento deste extenso dominio, apresentando divisões para os problemas em questão.

Uma dessas classificações, proposta por DeNardis e Raymond, segmenta a Governança da Internet em seis áreas funcionais:

\begin{enumerate}
  \item Controle de Recursos Críticos da Internet (CIRs);
  \item Estabelecimento de \textit{standards} da Internet;
  \item Coordenação de Acesso e Interligação;
  \item Governança de Cibersegurança;
  \item Intermediação de Informações;
  \item Direitos de Propriedade Intelectual baseados em Arquitetura;
\end{enumerate}

Estas áreas têm diferentes estratégias e modelos de admistração, cada um adaptado
para ser mais eficiente na sua área especifica. No entanto, todos os pontos 
convergem para um objetivo comum.

Examinando a trajetória da Internet ao longo do tempo, percebe-se que ela foi 
observada em diferentes contextos históricos e culturais, acompanhada de 
diversas expectativas e decepções.

Por conta disso, foi preciso mantar uma abordagem de flexivel face ao desenvolvimento
da arquitetura da Internet, permitindo que novos atores com diferentes ideias entrassem 
e modificassem a estrutura conforme desejado.





\section{O inicio da Internet}
\label{sec:beggining}
\section{Evolução da Governança da Internet}
\label{sec:governanceevo}
O termo "Governança da Internet" não está em nenhum artigo escolar até 1995. No entanto começaram a surgir pesquisadores com propostas em como aplicar 
Governança e leis na Internet.

O surgimento desta área de pesquisa converge para acontecimentos em meados do século 19, 
em que a Internet começou a emergir como uma massiva mídia
causada por 3 principais eventos:

\begin{enumerate}
  \item O desenvolvimento e adoção do protocologo WWW (1989-1993);
  \item[] Introduzindo a ideia de páginas web acessíveis a todos e permitiu a criação 
de uma rede de informações interligadas que poderiam ser acessadas por meio de navegadores
da \textit{web}.
  
  \item O surgimento de provedores de serviços de Internet (ISPs);
  \item[] Levou à popularização do acesso à Internet em domicilios e empresas;

  \item A publicação de \textit{web browsers} gratuitos para utilização (1991);
  \item[]Os \textit{web browsers} geralmente eram comercializados e a disponibilidade 
dos mesmos facilitou o acesso à Internet e abriu portas para novos individuos na sua
população.

  \item Privatização da Internet e a sua abertura para uso comercial 
pela \textit{US National Science Foundation} (1995); 
  \item[]Antigamente, a infraestrutura principal da Internet era predominanmente 
controlada e financiada pelo governo ou instituições acadêmicas, mas a decisão de 
permitir o uso comercial abriu caminho para o desenvolvimento do comércio eletrónico.
\end{enumerate}

Istou levou ao levantamento de questões em como deveriamos lidar com a populariação da Internet.
O artico sobre a Governança na Internet por Hardy em 1994 questiona-se se deviamos lidar com 
as leis, costumes e regras associadas a sua Governança de maneira diferente a como lidamos 
no mundo real e refere a criação de uma "\textit{multistakeholder collaboration}" para lidar 
com essas questões.


\subsection{A explosão dos anos 2000 e a ICANN}
A explosão da bolha da Internet em 2000 resultou na consolidação e reestruturação da Internet
e como podiamos visualiza-la como uma rede que connecta não só militares e académicos, mas 
que connecta o Mundo.

A formação da ICANN surgiu na ideia de criar uma instituição para a coordenação global dos dominios
e IP's da Internet.

A pesquisa e debates durante este periodo mudaram drasticamente para se seria ideal 
criar verdadeiras instituições governamentais para a Internet.

Em relação á ICANN houve uma enorme controversia, já que era uma corporação privada sem
fins-lucrativos pertencente aos Estados Unidos que estava a ganhar autoridade global 
dos nomes de dominio e IPs. Esta desconfiança foi resolvida através de contratos privados
feitos para lidar com questões de politica em termos de concorrencia no mercado
comercial de nomes de dominios e alocacao de enderecos IP.

\subsection{Da ICANN até aos dias de hoje}
De 2003 a 2009 a Governança da Internet tornou-se totalmente reconhecida como um dominio 
de Governança global e o tema expandiu além da ICANN mas foi na \textit{World Summit on the Information Society} (WSIS)
que mudou como viamos a pratica da Governança.

Consistiu em dois eventos patrocinados pela ONU para criar uma plataforma \textit{multistakeholder}
para abordar questões mundiais e expandiu o ponto da Governança da Internet o publico
da Web. 

Foi decidido então centralizar o serviço de DNS através da ICANN e foi implementado um modelo de
governança corporativa que envolve os setores privados e publicos como responsáveis pela
manutenção da Internet.

\subsection{Conclusão}

A evolução da governança da Internet ao longo das últimas décadas tem sido marcada por 
uma transformação significativa. Inicialmente, a Internet surgiu como uma plataforma 
para conectar acadêmicos e militares, mas rapidamente se expandiu para se tornar uma rede 
global, acessível a indivíduos e organizações de todo o mundo.

Em conclusão, a governança da Internet continua a ser um campo de estudo em constante 
evolução, refletindo a rápida transformação da própria Internet. É um tema de grande
 importância acadêmica e prática, pois influencia a forma como a Internet é utilizada e 
 regulamentada em todo o mundo.
\section{A Organização da Internet}
\label{sec:organization}
A Internet é constituida por um modelo complexo de multiplas camadas e elementos interconectados.
Uma visão geral da estrutura da Internet consiste em:

\subsection{Infraestrutura Fisica}
No coração da Internet está a sua estrutura fisica, composta por uma rede de cabos 
de fibra óptica, cabos submarinos, satélites e torres de comunicação sem fio.
Essa infraestrutura é repsonsável por transmitir os dados entre os diferentes pontos
da rede, formando o \textit{backbone} essencial da Internet.

\subsection{Backbone da Internet}
Consiste em uma rede de alta velocidade que interliga os principais pontos de troca 
de trafego em todo o mundo. Este é operado por empresas de telecomunicações e ISPs.

\subsection{Provedores de Serviços de Internet (ISPs)}
Os ISPs desempenham um papel fundamental na organização da Internet, fornecendo acesso à
rede para utilizadores e empresas. Estas conectam os dispositivos à Internet por meio de 
diferentes tecnologias.

\subsection{Protocologos}
No nivel mais fundamental da Internet estão uma série de protocolos de comunicação, como 
o TCP/IP, que padronizam a troca de dados entre os dispositivos conectados. Estes garantem
uma comunicação eficiente e confiável em toda a rede global.

\subsection{Endereços IP e DNS}
São essenciais para a identificação e localização de recursos na Internet. A ICANN desempenha 
um papel central na coordenação global desses recursos, garantindo que cada dispositivos
e serviço tenha um identificado único reconhecivel.

\subsection{Conclusão}
A organização da Internet é uma combinação complexa de infraestrutura física, 
protocolos de comunicação e entidades reguladoras. Essa estrutura descentralizada 
e distribuída permite que a Internet seja uma rede global resiliente, escalável 
e capaz de atender às crescentes demandas de uma sociedade em constante evolução. 
\section{A Internet em Portugal}
\label{sec:portugalcase}

Portugal connectou-se à Internet no Outuno de 1991, como resultado do projecto "Serviço IP 
da RCCN". Em uma altura em que os pioneiros portugueses na Internet ainda só estavam a 
começar a tomar contacto com a Internet, nos US já era possível trocar \textit{e-mails}
entre computadores através de ARPANET.

\subsection{Acesso da Internet aos portugueses}
Entre 1991 e 1995, para além das universidades e instituos de investigação portugueses 
apoiados pela FCCN, as instituições e individuos que tiveram acesso a serviços da Internet
fizeram-no via a PUUG (Portuguese Unix Users Group), ou pela INESC (Instituto de Engenharia
de Sistemas e Computadores) sendo sócios dessas associações sem fins lucrativos.

\subsection{Modernização dos portugueses na Internet}
O dominio DNS ".pt" foi definido em Setembro de 1991 pelo NIC (Network Information Center)
reconheceu a sua existência, posicionando Portugal no mapa da Internet Global, já que  
os servidores sob nosso controlo seriam vistos na Internet com o nome “.pt”.

Por volta dos anos 2000, o panorama do acesso à Internet ja estava totalmente mudado com a 
maioria dos acessos providenciada de forma pública por operadores de telecomunicações
nacionais e internacionais. Esse aumento do acesso público à Internet democratizou o 
conhecimento e abriu novas oportunidades para os portugueses.

\subsection{Dificuldades dos Portugueses}
Apesar dos avanços, os portugueses enfrentam desafios significativos na era digital, 
incluindo questões de segurança cibernética, privacidade dos dados e exclusão digital. 

\subsection{Conclusão}
A modernização dos portugueses na Internet representa uma jornada contínua de 
transformação e adaptação à era digital. À medida que Portugal avança rumo a um 
futuro digital, é crucial enfrentar os desafios e aproveitar as oportunidades 
oferecidas pela Internet, garantindo que todos os cidadãos possam colher os benefícios 
de uma sociedade digital inclusiva e próspera.

\section{O futuro da Internet}
À medida que avançamos para o futuro, a Internet continua a desempenhar um papel 
central em nossas vidas, impulsionando a inovação, conectando pessoas e transformando 
a sociedade. 

No entanto, à medida que exploramos as possibilidades do amanhã, enfrentamos uma série 
de desafios e questões complexas que moldarão o curso da Internet ao longo do seu tempo 
de vida.

\subsection{Desafios de Segurança e Privacidade}
À medida que a Internet se torna mais onipresente e integrada em nossas vidas,
vimos a enfrentar desafios crescentes relacionados à segurança cibernética e à 
privacidade dos dados. O aumento crescente de dispositivos conectados, 
da sofisticação de ciberataques e da coleta de dados pessoais levantam preocupações
sobre como proteger as nossas informações.

\subsection{Computação Quântica}
Um dos principais objetivos da humanidade é a realização da computação quântica,
uma tecnologia revolucionária que promete transformar como vemos a Internet e como 
processamos informações. 

Esta tem capacidade de processamento exponencialmente superiores às dos computadores clássicos
e tem o potencial de resolver problemas complexos e realizar cálculos que estão além do alcance 
da computação tradicional.

\subsection{Guerra da Inteligência Artificial}
Enquanto exploramos as fronteiras das Inteligências artificais, nos deperamos com 
uma nova competição e colaboração entre sistemas de IA.

Esta guerra, tanto no campo militar quanto comercial, levanta questões éticas, legais e de 
segurança sobre o uso e o desenvolvimento de IA's como por exemplo o controle das máquinas
sobre os humanos. 

\subsection{Conclusão}
À medida que nos aventuramos no futuro da Internet, é imperativo abordar esses desafios 
de forma inteligente e colaborativa. Ao enfrentar questões como computação quântica, 
guerra de inteligência artificial e segurança cibernética podemos moldar um futuro onde a 
Internet continue a ser uma força positiva para a humanidade, capacitando-nos a alcançar 
novos patamares de progresso, inovação e conectividade global. 

Estes são pontos criticos de debate e reflexão à medida que avançamos para um 
mundo cada vez mais auotmatizado e interconectado.
\section{Conclusão}
À medida que refletimos sobre o passado, examinamos o presente e imaginos o futuro 
da Internet, vemos uma historia repleta de desafios e promessas. 

Desde os primórdios da ARPANET, impulsionada pela necessidade de comunicação em tempos de 
incerteza global durante a Guerra Fria, até a globalização da Internet das nossas vidas 
é de certeza uma transformação extraordinária.

No entanto, à medida que avançamos para o futuro, enfrentamos desafios complexos que 
exigem soluções inovadoras e colaborativas. A segurança cibernética e a privacidade dos 
dados são questões prementes que exigem atenção e ação imediata.
Além disso, os avanços na computação quântica e a crescente competição na guerra de 
inteligência artificial trazem novas oportunidades e dilemas éticos que devem ser abordados 
de forma cuidadosa e responsável.

No centro de tudo isso, a governança da Internet desempenha um papel crucial na 
definição do futuro da rede global. Com múltiplas partes interessadas e áreas funcionais 
diversas, a governança da Internet é um campo em constante evolução, adaptando-se às 
necessidades e desafios do mundo digital em rápida transformação.

A Internet é algo que nos connecta sem nós termos noção disso, e para garantir que continua 
a ser uma força positiva para a humanidade, é imperativo permanecermos vigilantes 
e comprometidos com os valores fundamentais de uma Internet aberta, inclusiva e segura.

\nocite{mueller2020inventing}
\nocite{musiani2013network}
\nocite{lukasik2010arpanet}
\nocite{publicoPioneirosInternet}
\nocite{Tictank}

\bibliographystyle{plainnat}
\bibliography{references}
\end{document}